% =====================================================================
% APPENDICE — Indice (ceri e decisioni di design)
% =====================================================================
% Lo switch \appendix vive in main.tex prima dell'\input di questo file.

\chapter{Ceri e decisioni di design}\label{cap:ceri}

\noindent {\spacedlowsmallcaps{Questo \`e l'indice strutturale del libro}}. Ogni cosa che il libro assume senza dimostrarla --- ogni cosa che si \`e dovuta accendere come cero --- \`e qui. Separata da: ogni cosa che il libro \emph{decide} senza essere costretto a decidere cos\`i --- ogni scelta di design che ci sta come ci starebbe un'altra. Confondere le due cose \`e met\`a del lavoro mancato.

\section*{Ceri (assunzioni)}

\begin{center}
\renewcommand{\arraystretch}{1.4}
\begin{tabular}{p{0.5cm}p{5.2cm}p{3cm}p{4.8cm}}
\toprule
\textit{N.} & \textit{Assunzione} & \textit{Dove vive} & \textit{Perch\'e regge} \\
\midrule
1 & Esiste un metalinguaggio con i quattro giudizi & meta & \`e la lavagna su cui scrivi tutto il resto \\
\ref{cero:contesti} & Esistono i contesti e l'operazione $\vdash$ & meta & senza, niente ragionamento sotto ipotesi n\'e tipi dipendenti \\
\ref{cero:comp} & Le regole di computazione sono definitorie & meta & rispecchiano la struttura, non si provano dall'interno \\
\ref{cero:dipendenza} & Un termine pu\`o apparire dentro un tipo & scelta fondativa & \`e \emph{la} decisione che definisce MLTT \\
\midrule
(limite) & Coerenza del sistema & meta-meta & non dimostrabile dall'interno --- G\"odel \\
\bottomrule
\end{tabular}
\end{center}

\bigskip

\section*{Decisioni di design (scelte esplicite, non ceri)}

\begin{itemize}
\item \textbf{Eliminatore induttivo invece del solo case-split} per tipi ricorsivi --- $\indN$ vs $\mathsf{case\text{-}\mathbb{N}}$ (\S\ref{cap:bool} \emph{vs} \S\ref{cap:nat}). Potere dimostrativo: senza l'ipotesi induttiva nel passo, su $\mathbb{N}$ non si dimostra quasi niente. Si \emph{sarebbe potuto} scegliere solo case-split, e si avrebbe avuto un sistema diverso, pi\`u debole.

\item \textbf{Indice invece di parametro} per $n$ in $\mathrm{Vec}\,A\,n$ (\S\ref{cap:vec}). Permette a \texttt{cons} di cambiare la lunghezza producendo $\mathrm{Vec}\,A\,(\suc\,n)$ da $\mathrm{Vec}\,A\,n$. Se $n$ fosse parametro, sarebbe fisso per tutti i costruttori e \texttt{cons} non potrebbe esistere.

\item \textbf{Argomenti impliciti con \texttt{\{\}}} (\S\ref{cap:pi-azione}). Pura ergonomia, ma operativamente importante: \`e il riconoscimento sintattico che certi argomenti vivono solo nella passata 1 (type checking). L'argomento c'\`e sempre, sotto la sintassi.

\item \textbf{$\forall$ come zucchero per $\Pi$} (\S\ref{cap:forall}). Leggibilit\`a. Non \`e una regola nuova; \`e $\Pi$ con il permesso di omettere annotazioni di tipo quando Agda pu\`o inferirle.

\item \textbf{Parentesi raggruppate} \texttt{(a b c : A)} come zucchero per $\Pi$ annidati (\S\ref{cap:pi-azione}). Stessa famiglia di scelte: comodit\`a, sotto resta la regola di formazione unaria.
\end{itemize}

\bigskip

\noindent Ogni voce sopra \`e un punto in cui \emph{avresti potuto scegliere diversamente}. Saperlo \`e l'obiettivo: non confondere mai un'assunzione con un teorema, n\'e uno zucchero con una regola primitiva.

\bigskip
\bigskip

\section*{Dove vive ciascuna cosa --- la tavola finale}

\begin{center}
\renewcommand{\arraystretch}{1.25}
\begin{tabular}{p{6cm}p{3cm}p{4cm}}
\toprule
\textit{Cosa} & \textit{Dove vive} & \textit{Si dimostra?} \\
\midrule
$\tto, \ff, \zero, \suc\,n, \refl, \nil, \cons, \mathsf{inj}_1, \mathsf{inj}_2, \mathsf{fz}, \mathsf{fs}$ & linguaggio & s\`i, sono termini (introdotti dalle regole I) \\
$\mathbb{B}, \mathbb{N}, \Pi, \Sigma, A \uplus B, \_\identita\_, \mathrm{Fin}\,n, \mathrm{Vec}\,A\,n, \bot$ & linguaggio & s\`i, sono tipi (introdotti dalle regole F) \\
$\indB, \indN, \indS, \mathsf{ind}\text{-}\uplus, \Jelim$, applicazione di $\Pi$ & metalinguaggio & no, sono regole (E) \\
Regole di computazione (Cero~\ref{cero:comp}) & metalinguaggio & no, sono definitorie (C) \\
$\definitorio$ definitorio (giudizio) & metalinguaggio & no, \`e un'affermazione \\
$\identita$ proposizionale (tipo $\_\identita\_$) & linguaggio & s\`i, si abita con $\refl$ \\
Coerenza del sistema & meta-meta & non dall'interno (G\"odel) \\
\bottomrule
\end{tabular}
\end{center}

\bigskip

\noindent Pi\`u si scende verso le fondamenta, pi\`u il terreno diventa silenzioso. A un certo punto c'\`e sempre un cero alla madonna --- e va bene cos\`i, \emph{purch\'e sia esplicito}.

\bigskip

\noindent Lo svolgimento concreto del libro di Stump --- capitolo per capitolo, F-I-E-C alla mano --- \`e in Parte~IV (\S\ref{cap:stump-1} e seguenti).
